\section*{Chapter \ref{ch:dv:parametrising-the-surface} --- Parametrising the Surface}

\begin{solution}{exo:dv:parametrising-the-surface:1}
$w(0)=a+b\bigl(-\rho m+\sqrt{m^2+\varsigma^2}\bigr)=0.00670$; the implied volatility is
$\sqrt{0.00670/(91/365)}=16.4\%$.
\end{solution}

\begin{solution}{exo:dv:parametrising-the-surface:2}
Left $b(1-\rho)=0.062$, right $b(1+\rho)=0.014$. From $\tilde q=1/(2\beta_L)+\beta_L/8-1/2$ with
$\beta_L=0.062$: $\tilde q=7.6$, so negative moments $\E S_T^{-q}$ are finite up to about $q=7.6$; far
from the bound 2, which would mean no finite negative moment.
\end{solution}

\begin{solution}{exo:dv:parametrising-the-surface:3}
$0.5^2\times12/365-0.3^2\times12/365=0.00526$; standard deviation 7.3\%; expected absolute move
$7.3\%\times\sqrt{2/\pi}=5.8\%$.
\end{solution}

\begin{solution}{exo:dv:parametrising-the-surface:4}
The event adds variance at one date, not at a rate. Linear interpolation spreads it over the interval,
so an option expiring the day before the release is charged part of the event's variance it will never
see: on the chapter's numbers a 7-day option would be priced at 49.5\% instead of 32\%.
\end{solution}

\begin{solution}{exo:dv:parametrising-the-surface:5}
$\varphi=1.095\times0.00183^{-0.478}$; $\theta\varphi(1+|\rho|)=0.064<4$ and
$\theta\varphi^2(1+|\rho|)=1.42\le4$: both hold.
\end{solution}

\begin{solution}{exo:dv:parametrising-the-surface:6}
$\sigma_c=0.25\sqrt{(5/252)/(7/365)}=25.43\%$.
\end{solution}

\begin{solution}{exo:dv:parametrising-the-surface:7}
On these quotes both fits lie inside all 27 bands and their wings are within $2\times10^{-4}$ of each
other (0.0623 and 0.0140 against 0.0621 and 0.0142). The band fit matters when quotes are wide or
lopsided, and in the wings, where mids are least reliable.
\end{solution}

\begin{solution}{exo:dv:parametrising-the-surface:8}
The mids are not the fair values: they carry noise of a fraction of the spread. Passing through them
reproduces the noise (the spline is 0.17 points from the true smile at worst, SVI 0.08), puts it into
the density and the \omterm{def:dv:greeks-and-the-hedging-pnl:greeks}{Greeks}, and says nothing about the wings. The right error measure is distance to the
bid--ask band, which both satisfy, and stability of the surface from one refit to the next.
\end{solution}

\begin{solution}{pb:dv:parametrising-the-surface:1}
\textbf{1.} $0.32^2\times4/365=0.001122$; $0.561^2\times11/365=0.009485$.
\textbf{2.} $0.32^2\times11/365=0.003086$.
\textbf{3.} $0.009485-0.003086=0.0064$.
\textbf{4.} 8.0\% and 6.4\%.
\textbf{5.} The 4-day expiry is the only one that ends before the release; later expiries all contain
the event.
\textbf{6.} 32.0\%.
\textbf{7.} 32.1\%.
\textbf{8.} A second event before those expiries, or a market that expects higher diffusive volatility
after the release.
\textbf{9.} About 32\%: the \omterm{def:dv:parametrising-the-surface:event}{event variance} leaves the expiry's total variance once the release is past.
\textbf{10.} The straddle was priced with the event's variance; a 6\% move is below the 8\% standard
deviation priced, and the volatility crush after the release takes the time value with it.
\textbf{11.} Linearly in total variance on a clock that puts the event's variance on day 7 and the
diffusive variance on the other trading days.
\textbf{12.} The expiries just after the release: a jump makes the smile more curved at short expiries,
and more so the larger the \omterm{def:dv:parametrising-the-surface:event}{event variance} relative to the diffusive variance.
\textbf{13.} Give the weekend days little or no variance in \omterm{def:dv:parametrising-the-surface:business}{business time}.
\textbf{14.} The diffusive 32\%.
\textbf{15.} The slice's level falls sharply and the curvature that the jump added near the money goes;
the refit's parameters jump accordingly, which is why parameter changes on event dates are expected, not
alarms.
\textbf{16.} No: it is the move priced under the pricing measure, including a premium for bearing the
event risk.
\textbf{17.} Compare the implied expected absolute move with the average absolute move on the share's past
releases, measured the same way (close before to close after).
\textbf{18.} Sell the 11-day straddle and buy the 4-day or a later expiry to hedge the diffusive
volatility; the risk is a move larger than priced, a loss of gamma on the release day.
\textbf{19.} A standard deviation of 8.0\% and an expected absolute move of 6.4\%.
\textbf{20.} Because an event adds variance on one date, and interpolation in time would spread it over
expiries that end before it.
\end{solution}

\begin{solution}{iq:dv:parametrising-the-surface:1}
$w(k)=a+b(\rho(k-m)+\sqrt{(k-m)^2+\varsigma^2})$: $a$ raises the whole smile, $b$ opens the angle between
the wings, $\rho$ rotates it (skew), $m$ shifts it left or right, $\varsigma$ rounds the vertex.
\iqlookfor{the formula and the geometric role of each parameter.}
\end{solution}

\begin{solution}{iq:dv:parametrising-the-surface:2}
Compute total variances; the \omterm{def:dv:parametrising-the-surface:event}{event variance} is the later expiry's total variance minus the earlier
expiry's variance rate times the later expiry's time; its square root is the standard deviation of the
move, and $\sqrt{2/\pi}$ times it the expected absolute move.
\iqlookfor{working in total variance, not volatility.}
\end{solution}

\begin{solution}{iq:dv:parametrising-the-surface:3}
Lee's moment formula: the slope of total variance in the wing is $2-4(\sqrt{p^2+p}-p)$, where $p$ counts
the finite moments of the underlying; it is at most 2, attained when no moment beyond the first is
finite. A steeper wing would give call prices that violate arbitrage bounds at extreme strikes.
\iqlookfor{the bound, and its link to moments.}
\end{solution}

\begin{solution}{iq:dv:parametrising-the-surface:4}
Reduce the dimension: for fixed $(m,\varsigma)$ SVI is linear in three parameters, solved by least squares;
search the remaining two from yesterday's (or the last second's) values; impose Lee's bound and the
admissible domain; minimise distance outside the bands with weights by width; check arbitrage after each
fit and fall back to the previous surface if a check fails.
\iqlookfor{dimension reduction, warm starts, bands, checks.}
\end{solution}

\begin{solution}{iq:dv:parametrising-the-surface:5}
Friday's marks include two days of calendar \omterm{def:dv:greeks-and-the-hedging-pnl:greeks}{theta} that the market does not charge; on Monday the
options are re-marked higher than the model expects, and the desk books an apparent \omterm{def:dv:greeks-and-the-hedging-pnl:greeks}{vega} gain that is
the weekend \omterm{def:dv:greeks-and-the-hedging-pnl:greeks}{theta} coming back. Move to \omterm{def:dv:parametrising-the-surface:business}{business time}, with weekend weights.
\iqlookfor{the clock as a modelling choice.}
\end{solution}

\begin{solution}{iq:dv:parametrising-the-surface:6}
SSVI when the surface must be free of arbitrage by construction and stable (risk systems, exotics
pricing, sparse or illiquid expiries), and to extrapolate in expiry; independent slices when each expiry
must match its quotes closely (market making). SSVI gives up the fit of short, steep smiles, since one
$\rho$ and one curvature law span all expiries.
\iqlookfor{guarantees against flexibility.}
\end{solution}
